\begin{document}
\maketitle
% 原文の引用番号を維持するため，原文の文献順を明示する．
\nocite{Guth:1980zm,Linde:1981mu,Albrecht:1982wi,Hubble,abc,DickeCMB,WhitePaper,Dodelson,Mukhanov,WeinbergCosmology,WMAP5,unionsn,BassettReview,PaoloEFT,WeinbergEFT,BMReview,Shafi:2006cs,LythRiotto,WandsMulti,Bardeen:1980kt,BST,Wands:2000dp,Acquaviva02,malda,BD,Jacobson:2003vx,Lyth,WayneHu,CMBFAST,CAMB,NGreview,Kamionkowski:1996ks,Zaldarriaga:1996xe,NewmanPenrose,spergel/zaldarriaga:1997,Dodelson:2003ip,EPIC,BBKS,EisensteinHu,Steinhardt:2002ih,Steinhardt:2001st,Buchbinder:2007ad,Smith:2009jr,Nflation,MSW,Silverstein:2008sg,KKLMMT,Baumann:2007ah,Baumann:2008kq,Conlon:2005jm,BlancoPillado:2006he,BauZal,Babich_etal_04,BMRmulti,Bernardeau:2002jy,Bernardeau:2002jf,Sasaki:2008uc,Naruko:2008sq,Byrnes:2008wi,Byrnes:2006fr,Langlois:2008vk,Valiviita:2008zb,Assadullahi:2007uw,Valiviita:2006mz,Vernizzi:2006ve,Allen:2005ye,Linde:1996gt,Lyth_etal03,Dvali:2003em,Kofman:2003nx,Creminelli:2007aq,Koyama_etal07,Buchbinder_etal07,Lehners:2007wc,Lehners:2008my,Koyama:2007ag,Koyama:2007mg,Alishahiha:2004eh,GhostInflation,Senatore:2004rj,Creminelli:2003iq,Chen:2006nt,Gordon:2000hv,Lyth:2002my,Weinberg04b,Komatsu:2009kd,Schwinger:1960qe,Calzetta:1986ey,Jordan:1986ug,Weinberg:2005vy,komatsu01,Fergusson:2008ra,DBI,Holman:2007na,Creminelli:2004yq,Cheung:2007sv,Senatore:2009gt,BMR1,BMR2,Senatore:2008wk,Dalal:2007cu,Slosar:2008hx,Copeland:1994vg,PolchinskiOne,Polchinski,KKLT,Lyth:1996im,Linde:1983gd,Vafa:2005ui,Ooguri:2006in,Adams:2006sv,Hertzberg:2007ke,Douglas:2006es,KS,KT,GKP,Baumann:2006cd,Baumann:2006th,Baumann:2007np,Krause:2007jk,Nakahara,BaumannProgress,Firouzjahi:2005dh,ShanderaTye,Freese:1990rb,Banks:2003sx,Svrcek:2006yi,Easther:2005zr,Liddle:1998jc,Kallosh:2007cc,Grimm:2007hs,Hertzberg:2007wc,Flauger:2008ad,Caviezel:2008tf,Sarangi,CMP,Burgess:2007pz,Kallosh:2007ig,Linde:2005dd,McAllister:2007bg,Cline:2006hu,Flauger:2009ab,:2006uk,CMBPolWorkshop,Kosowsky:2004sw,Ruhl:2004kv,Taylor:2006jw,Samtleben:2008rb,Yoon:2006jc,Oxley:2005dg,MacTavish:2007kh,Bock:2009xw,Vilenkin:1983xq,Linde:1986fd,Linde:1993nz,Linde:1993xx,GarciaBellido:1993wn,Vilenkin:1994ua,Garriga:2005av,Bousso:2006ev,Linde:2007nm,Bousso:2008hz,DeSimone:2008if,Garriga:2008ks,Winitzki:2008jp,Linde:2008xf,Malik:2008im,LiciaStats,ADM,Seery:2007we,Adshead:2009cb}
\begin{center}
{\small 日本語訳：\href{https://arxiv.org/abs/0907.5424v2}{arXiv:0907.5424v2}\\
原題：TASI Lectures on Inflation}
\end{center}
\begin{abstract}
5回の連続講義を通じて，インフレーション宇宙論を概観する．まずFriedmann--Robertson--Walker（FRW）宇宙論における初期条件の問題を説明し，宇宙初期の加速膨張期であるインフレーションが，これらの問題をどのように解決するかを述べる．次に，インフレーションが微視的な量子揺らぎを宇宙の構造形成の巨視的な種へと変換する仕組みを説明する．スカラー揺らぎとテンソル揺らぎの原始スペクトルを求める有名な計算を，詳細に提示する．そのうえで，宇宙マイクロ波背景放射の揺らぎの観測からインフレーション期の情報を抽出するという逆問題を定式化する．インフレーションを支持する現在の観測的証拠と，今後の検証の可能性を論じる．最後に，弦理論におけるインフレーションを解説し，インフレーションを基礎物理と結びつける際の課題を概観する．
\end{abstract}

\paragraph{第1講：インフレーションの古典力学}
本講の目的は，インフレーション宇宙論の古典力学を第一原理から導入することである．基本的なFRW宇宙論を簡単に復習した後，従来のビッグバン理論が初期条件の問題をもたらすことを示す．すなわち，私たちが知るような宇宙は，非常に特殊で微調整された初期条件からしか生じない．次に，インフレーション，すなわち宇宙初期の加速膨張期がこの初期条件の問題を解決し，一般的な初期条件から私たちの宇宙が生じることを可能にする仕組みを説明する．インフレーションに必要な条件を述べ，インフレーションが時空の因果構造を変えることによって，ビッグバンの難問をどのように解決するかを説明する．最後に，インフレーション的膨張の物理的起源を論じて本講を終える．

\paragraph{第2講：インフレーション中の量子揺らぎ}
本講では，インフレーション中の量子揺らぎが生成する原始揺らぎスペクトルの有名な計算を復習する．計算をすべて詳しく示し，「ごまかし」や近似をできるだけ避ける．宇宙論的摂動論の基礎的な側面を簡単に復習した後，まず，インフレーションが微視的な量子揺らぎを宇宙の構造形成の巨視的な種へ変換する基本的な仕組みを定性的にまとめる．次いで，曲がった時空の場の量子論への教育的な導入として，単純な調和振動子の量子化を復習する．真空が最低エネルギー状態であることを要求すると，一意な真空状態が選ばれる点を強調する．その後，インフレーションに対する対応する計算を行う．スカラー揺らぎとテンソル揺らぎの両方のパワースペクトルを計算する．

\paragraph{第3講：観測との接点}
本講では，宇宙マイクロ波背景放射と大規模構造の観測から，インフレーションによる摂動スペクトルの情報を抽出するという逆問題を説明する．前講で計算したゲージ不変なスカラーおよびテンソルのパワースペクトルと，観測されるCMB異方性および銀河パワースペクトルの間の正確な関係を定義する．原始揺らぎと後期宇宙の観測量を関係づける伝達関数を示す．これらの結果を用いて，インフレーションを支持する現在の観測的証拠を論じる．最後に，今後インフレーションを検証する機会を示す．

\paragraph{第4講：原始非ガウス性}
本講では，原始非ガウス性の研究における主要な理論的結果をまとめる．結果の大部分は証明なしに述べるが，インフレーションの根本的な物理的起源を制約するうえでの意義を説明する．非ガウス統計の基本的な診断量としてバイスペクトルを導入した後，その運動量依存性がインフレーションの作用を探る強力な手段になることを示す．大きな非ガウス性は，インフレーション中にインフラトンの相互作用が重要である場合にのみ生じうる．したがって，単一場スローロール・インフレーションでは，非ガウス性は観測できないほど小さいと予言される一方，複数の場，高階微分相互作用，または非標準的な初期状態をもつ模型では，有意な大きさになりうる．最後に，原始非ガウス性を検出または制約する観測的な見通しを論じて本講を終える．

\paragraph{第5講：弦理論におけるインフレーション}
この連続講義を，やや発展的な話題である弦理論におけるインフレーションの議論で締めくくる．Liam McAllisterとの最近の総説に基づき，この主題を教育的に概観する．本講の中心的な主題は，インフレーションのPlanckスケールの物理に対する敏感さである．これこそが，弦理論でインフレーションを実現する主な動機であると同時に，中心的な理論課題でもあると論じる．これらの問題を，弦理論における2つのインフレーション・シナリオ，すなわちワープしたDブレーン・インフレーションとアクシオン・モノドロミー・インフレーションの事例を通じて説明する．最後に，理論と観測の両面で今後進展が期待される方向を示す．

\vspace{1em}
\noindent{\footnotesize 電子メール：\texttt{dbaumann@physics.harvard.edu}}
\setcounter{footnote}{0}
\clearpage
\hypertarget{toc}{}
\tableofcontents
\clearpage
